\chapter{Feature Dictionary}

This appendix documents the structured variables used in the primary triage-based clinical and multimodal models. The active feature policy is intentionally restricted to variables that are available at or before the imaging time ($t_0$), in order to preserve temporal validity and reduce the risk of information leakage from downstream clinical decisions.

The primary structured feature set is derived from emergency department triage data linked to the chest radiograph cohort. It consists of four groups of variables: numeric physiological variables, categorical context variables, explicit missingness indicators for numeric fields, and binary view-position flags inherited from the imaging pipeline. The same core triage feature set is used in the headline clinical logistic regression model, the headline clinical XGBoost model, and the structured branch of the headline multimodal model.

Feature construction is implemented through the triage linkage and feature-building pipeline, primarily in \texttt{src/data/link\_cxr\_to\_triage.py}, \texttt{src/data/build\_triage\_features.py}, and \texttt{src/data/build\_triage\_model\_table.py}. The active feature policy is also described in \texttt{docs/triage\_feature\_policy.md}. Missingness counts are recorded in artifact reports, but precomputed missingness percentages were not found as direct stored outputs in the repository. Likewise, units are only partially explicit in code comments; beyond temperature, exact units are generally not encoded directly in the stored feature tables and are therefore treated here as \emph{not explicitly documented in code} unless otherwise stated.

\section{Included Features}

The features included in the primary triage-based models are grouped as follows.

\subsection*{Numeric Variables}
\begin{itemize}
    \item \texttt{temperature}
    \item \texttt{heartrate}
    \item \texttt{resprate}
    \item \texttt{o2sat}
    \item \texttt{sbp}
    \item \texttt{dbp}
    \item \texttt{pain}
    \item \texttt{acuity}
\end{itemize}

\subsection*{Missingness Indicators}
\begin{itemize}
    \item \texttt{temperature\_missing}
    \item \texttt{heartrate\_missing}
    \item \texttt{resprate\_missing}
    \item \texttt{o2sat\_missing}
    \item \texttt{sbp\_missing}
    \item \texttt{dbp\_missing}
    \item \texttt{pain\_missing}
    \item \texttt{acuity\_missing}
\end{itemize}

\subsection*{Binary Flags}
\begin{itemize}
    \item \texttt{is\_pa}
    \item \texttt{is\_ap}
\end{itemize}

\subsection*{Categorical Variables}
\begin{itemize}
    \item \texttt{gender}
    \item \texttt{race}
    \item \texttt{arrival\_transport}
\end{itemize}

\section{Preprocessing and Feature Policy}

Numeric variables are coerced to numeric type, with implausible values clipped for several vital-sign fields during triage feature construction. In the logistic regression and multimodal pipelines, these variables are subsequently median-imputed and standardized. In the XGBoost pipeline, numeric values are used after coercion without an explicit scaling stage. Categorical variables are normalized to strings, trimmed, and mapped to \texttt{UNKNOWN} when missing or empty. In the logistic regression and multimodal pipelines, they are imputed with a constant \texttt{UNKNOWN} token and one-hot encoded; in XGBoost, they are cast to categorical type for native handling. Missingness indicators are explicitly created during feature construction and then filled to one when absent during model preparation.

Several plausibility filters are explicitly implemented in \texttt{src/data/build\_triage\_features.py}. Temperature is clipped to the range 95.0 to 105.8, with a code comment indicating Fahrenheit. Heart rate is clipped to 30--220, respiratory rate to 5--60, oxygen saturation to 50--100, systolic blood pressure to 60--250, and diastolic blood pressure to 30--150. No explicit clipping is implemented for \texttt{pain} or \texttt{acuity} in the same feature-construction stage.

The two binary flags, \texttt{is\_pa} and \texttt{is\_ap}, are not native triage variables but propagated imaging-view indicators inherited from the radiograph selection pipeline. They are included as structured inputs in all primary model families because view position may affect radiographic appearance and therefore influence prediction behavior.

\section{Excluded or Leakage-Sensitive Fields}

Two fields are explicitly notable in the repository but are not part of the active primary baseline feature set:

\begin{itemize}
    \item \texttt{disposition}: explicitly removed or marked as leakage-prone in both documentation and feature-building code because it may reflect downstream post-\(t_0\) clinical decisions.
    \item \texttt{chiefcomplaint}: present at intermediate stages of feature building, but explicitly excluded from the baseline clinical and multimodal models.
\end{itemize}

This exclusion policy is important for the thesis because it demonstrates that the structured branch was designed around time-safe variables rather than convenience or maximal feature inclusion.

\section{Detailed Feature Table}

\renewcommand{\arraystretch}{1.15}
\small
\begin{longtable}{p{2.6cm} p{2.2cm} p{2.3cm} p{6.8cm}}
\toprule
\textbf{Feature} & \textbf{Type} & \textbf{Used In} & \textbf{Preprocessing / Notes} \\
\midrule
\endfirsthead

\toprule
\textbf{Feature} & \textbf{Type} & \textbf{Used In} & \textbf{Preprocessing / Notes} \\
\midrule
\endhead

\texttt{temperature} & Numeric & All primary models & Numeric coercion; clipped to 95.0--105.8; median imputation + standardization in logistic/multimodal; raw coerced numeric in XGBoost. Code comment indicates Fahrenheit. \\
\texttt{heartrate} & Numeric & All primary models & Numeric coercion; clipped to 30--220; median imputation + standardization in logistic/multimodal; raw coerced numeric in XGBoost. \\
\texttt{resprate} & Numeric & All primary models & Numeric coercion; clipped to 5--60; median imputation + standardization in logistic/multimodal; raw coerced numeric in XGBoost. \\
\texttt{o2sat} & Numeric & All primary models & Numeric coercion; clipped to 50--100; median imputation + standardization in logistic/multimodal; raw coerced numeric in XGBoost. \\
\texttt{sbp} & Numeric & All primary models & Numeric coercion; clipped to 60--250; median imputation + standardization in logistic/multimodal; raw coerced numeric in XGBoost. \\
\texttt{dbp} & Numeric & All primary models & Numeric coercion; clipped to 30--150; median imputation + standardization in logistic/multimodal; raw coerced numeric in XGBoost. \\
\texttt{pain} & Numeric & All primary models & Numeric coercion; no explicit clipping found; median imputation + standardization in logistic/multimodal; raw coerced numeric in XGBoost. \\
\texttt{acuity} & Numeric & All primary models & Numeric coercion; no explicit clipping found; median imputation + standardization in logistic/multimodal; raw coerced numeric in XGBoost. \\
\texttt{temperature\_missing} & Missingness indicator & All primary models & Derived as missingness flag; filled to 1 and cast to integer in model preparation. \\
\texttt{heartrate\_missing} & Missingness indicator & All primary models & Derived as missingness flag; filled to 1 and cast to integer in model preparation. \\
\texttt{resprate\_missing} & Missingness indicator & All primary models & Derived as missingness flag; filled to 1 and cast to integer in model preparation. \\
\texttt{o2sat\_missing} & Missingness indicator & All primary models & Derived as missingness flag; filled to 1 and cast to integer in model preparation. \\
\texttt{sbp\_missing} & Missingness indicator & All primary models & Derived as missingness flag; filled to 1 and cast to integer in model preparation. \\
\texttt{dbp\_missing} & Missingness indicator & All primary models & Derived as missingness flag; filled to 1 and cast to integer in model preparation. \\
\texttt{pain\_missing} & Missingness indicator & All primary models & Derived as missingness flag; filled to 1 and cast to integer in model preparation. \\
\texttt{acuity\_missing} & Missingness indicator & All primary models & Derived as missingness flag; filled to 1 and cast to integer in model preparation. \\
\texttt{is\_pa} & Binary flag & All primary models & Propagated imaging-view flag from the frontal-image selection pipeline; filled to 0 when absent; treated as numeric input. \\
\texttt{is\_ap} & Binary flag & All primary models & Propagated imaging-view flag from the frontal-image selection pipeline; filled to 0 when absent; treated as numeric input. \\
\texttt{gender} & Categorical & All primary models & String normalization; missing or empty mapped to \texttt{UNKNOWN}; one-hot encoded in logistic/multimodal; categorical dtype in XGBoost. \\
\texttt{race} & Categorical & All primary models & String normalization; missing or empty mapped to \texttt{UNKNOWN}; one-hot encoded in logistic/multimodal; categorical dtype in XGBoost. \\
\texttt{arrival\_transport} & Categorical & All primary models & String normalization; missing or empty mapped to \texttt{UNKNOWN}; one-hot encoded in logistic/multimodal; categorical dtype in XGBoost. \\
\bottomrule
\end{longtable}
\normalsize

\section{Feature Provenance and Time-Safety Checks}

The primary feature set documented here is identical across the clinical logistic regression baseline, the clinical XGBoost baseline, and the structured branch of the multimodal model. This consistency is important because it ensures that performance differences among those systems arise from model family and modality interactions rather than from changes in the structured input definition.

Feature provenance and time-safety are supported by the following implementation components:

\begin{itemize}
    \item \texttt{src/data/build\_triage\_features.py}: constructs cleaned triage variables, plausibility filtering, and missingness indicators.
    \item \texttt{src/data/build\_triage\_model\_table.py}: defines the structured model table used by downstream training scripts.
    \item \texttt{docs/triage\_feature\_policy.md}: records the intended time-safe baseline policy and leakage-sensitive exclusions.
\end{itemize}

No precomputed feature-level missingness percentages were found as direct stored values in the repository, although missingness counts are available in artifact reports. Similarly, except for the temperature code comment, units are generally not stored as explicit metadata in the feature tables. These omissions do not affect the training pipeline itself, but they do limit the level of direct variable metadata available for reporting.

\section{Artifact References}

Relevant artifact and source locations include:

\begin{itemize}
    \item \texttt{src/models/clinical\_baseline.py}
    \item \texttt{src/models/clinical\_xgb.py}
    \item \texttt{src/training/train\_multimodal\_pneumonia.py}
    \item \texttt{src/data/build\_triage\_features.py}
    \item \texttt{src/data/build\_triage\_model\_table.py}
    \item \texttt{docs/triage\_feature\_policy.md}
    \item \texttt{artifacts/manifests/cxr\_ed\_triage\_features\_report.json}
    \item \texttt{artifacts/manifests/cxr\_ed\_triage\_model\_table\_report.json}
\end{itemize}